% New native coauthor theorem; inherited sources are unchanged.
% Dependencies: physical endpoint first variation, 281--284, 317, 322--324, 338.
\section{Exponential payment of interior section isoroof charges}
\label{sec:iso-section}

The critical incoming fans and the regular constant-roof side arcs
are different terms in the exact signed derivative. We now pay the
second term on every original interior section side. The proof does
not assert that a constant-roof arc is absent. It uses the actual
two-half optical balance to confine its endpoint momenta to a thin
angular sector, and a truncated interval-overlay estimate pays the
complementary derivative regime.

Retain $R\in[9/20,47/100]$, the physical collision count $m$, and all
canonical-source constants and the qualified enlarged curve class
of Section~\ref{sec:cs-canonical}. Put
\begin{equation}\label{eq:iso-constants}
 c_\circ=\sqrt{19}/10,\quad \kappa_-=100/47,\quad U=170/9,
 \quad Q=\frac{U}{c_\circ\kappa_-}=\frac{799}{9\sqrt{19}},
 \quad F_{\rm cs}=\frac{\Lambda_\circ}{4\pi c_*R_-^2}.
\end{equation}
The original endpoint section tests lie inside $|p_0|,|p_m|<5/6$,
so the retained cap $|p_0|,|p_m|\le9/10$ covers every coefficient
estimated below. It is not an added conditioning event. On this cap
the one-copy regular roof--angle charge is
\[
 A\,|d\theta|,\qquad
 \theta=\arg(-p_0,p_m),\qquad
 A=\frac{u_x}{4\pi c_*h_x}\le F_{\rm cs}u_x,
 \qquad h_x=R^2 e^{\mathsf T}G^{-1}e.
\]
These are the exact physical quantities of
Corollary~\ref{cor:bf-wronskian}, with one conditioning factor.

\subsection{The endpoint balance and the angular image}

Fix $1\le i\le m-1$ and an original canonical side.
Use $v=\alpha_i$ on $p_i=b$, or $v=\varphi_i$ on
$\alpha_i=a$. In either case the two true endpoint images are
unstable graphs in their respective physical time directions.
Write $V_-$ for the reversed graph slope and $V_+$ for the
forward slope in $(s,\varphi)$ coordinates. Thus
\[
 \frac{d(-\varphi_0)}{d(R\alpha_0)}=V_-,
 \qquad \frac{d\varphi_m}{d(R\alpha_m)}=V_+,
 \qquad \kappa_-\le V_-,V_+\le U.
\]
The source and all contacts are actual physical states. No segment
has been added to the roof.

\begin{lemma}[True two-half first variation and a single opposition quadrant]
\label{lem:iso-balance}
On every connected actual two-sided canonical overlap, put
$c_j=\sqrt{1-p_j^2}$ and $K_-=c_0V_-$, $K_+=c_mV_+$.
Then
\begin{equation}\label{eq:iso-first-variation}
 \frac{dF}{dv}=
       \frac{p_0}{K_-}\frac{dp_0}{dv}
       +\frac{p_m}{K_+}\frac{dp_m}{dv},
 \qquad c_\circ\kappa_-\le K_\pm\le U.
\end{equation}
Let $S$ be any measurable subset of this overlap on which
$F_v=0$ and $(p_0,p_m)\ne(0,0)$.
The momentum components are both nonzero on $S$. After fixed
sign changes making both components strictly increasing with $v$,
all of $S$ lies in a single opposition quadrant.
The Gauss angle is strictly monotone on that entire quadrant
interval. Consequently the total angular image of $S$ has length
at most $\pi/2$, without a component count.

Designate one half as the short half and the other as the long half,
breaking an equality by a fixed rule. Let
\[
 a(v)=|p_{\rm short}'(v)|^{-1},\qquad
 b(v)=|p_{\rm long}'(v)|^{-1}.
\]
At every point of $S$,
\begin{equation}\label{eq:iso-balance-ratio}
 \frac{|p_{\rm long}|}{|p_{\rm short}|}
       =\frac{bK_{\rm long}}{aK_{\rm short}}
       \le Q\frac ba.
\end{equation}
For any $\eta>0$ and $\epsilon\ge\sup b$, the angular image of
$S\cap\{a>\eta\}$ has length at most
$\arctan(Q\epsilon/\eta)\le Q\epsilon/\eta$.
\end{lemma}
\begin{proof}
The reversed endpoint graph has
$p_0'=-Rc_0V_-\alpha_0'$, whereas the forward graph has
$p_m'=Rc_mV_+\alpha_m'$. Insert both equalities in the true
endpoint first variation
\[
 F_v=-Rp_0\alpha_0'+Rp_m\alpha_m'.
\]
This proves \eqref{eq:iso-first-variation}, including the plus
sign in its backward term. The endpoint cap and the verified
unstable cone give its bounds for $K_\pm$.

Each endpoint momentum derivative is nonzero and has a fixed sign
on the actual canonical partition interval. If one endpoint
momentum vanishes at a point of $S$, the identity forces the other
to vanish, contrary to its definition. After the fixed sign changes
write the components as $X,Y$, with $X',Y'>0$. Stationarity gives
\[
             XX'/K_X+YY'/K_Y=0,
\]
so $XY<0$. An increasing planar curve can visit at most one of
the two opposition quadrants: a visit to both would force one
of its components to decrease across zero. Its intersection with
that quadrant is an interval. There
\[
 \frac{d}{dv}\frac{-Y}{X}
       =\frac{-Y'X+YX'}{X^2}
\]
has a fixed nonzero sign, positive for $X<0<Y$ and negative for
$Y<0<X$. Hence its angular coordinate is strictly monotone
throughout that interval. Fixed coordinate sign changes and
interchanging the two endpoints preserve absolute angular lengths.
This proves the assertion for the actual Gauss angle
$\arg(-p_0,p_m)$, including any disconnected selected subset.

Taking absolute values in the stationarity identity gives
\eqref{eq:iso-balance-ratio}. If $a>\eta$ and $b\le\epsilon$,
the selected directions lie in the single angular wedge with
$|p_{\rm long}|/|p_{\rm short}|\le Q\epsilon/\eta$ next to
the short-momentum axis. Since the angle on the whole opposition
interval is one-to-one, all selected components together have
angular variation at most the width of that one wedge.
No constancy, derivative bound, or finite-component property
is required of a selector.
\end{proof}

The lemma covers every regular constant-roof side arc by taking
$S$ to be its union with the other such arcs on the overlap.
It also bounds the complete regular stationary-roof subsource.
There is no assumption that a pulled-back section side is a straight
line in endpoint coordinates or that its roof restriction is convex.

\subsection{A truncated optical overlay pays both derivative regimes}

Write $H=H_{\rm hor}$ for a horizontal source and
$H=H_{\rm ver}$ for a vertical source. Put $E=E_{\rm cs}=47/6$,
$\sigma=47/59$, and
\[
 L=\max(i,m-i),\qquad s=\min(i,m-i),\qquad
 \epsilon=E\sigma^{L-1},\qquad A_s=E\sigma^{s-1}.
\]
Let $a_I$ and $b_J$ be the full interval suprema of $a$ and $b$
on the short and long half partitions. Theorem~\ref{thm:cs-half-sup}
gives, on the original common canonical source,
\begin{equation}\label{eq:iso-halves}
 \sum_Ia_I\le H,\quad \sum_Jb_J\le H,\quad
 a_I\le A_s\le E,\quad b_J\le\epsilon\le E.
\end{equation}
The endpoint cap leaves at most one subinterval of each one-sided
interval and does not increase these bounds. Put
\[
 C_{\rm side}=
 \begin{cases}
 C_0\sqrt{1-b^2},&\text{on }p_i=b,\\
 1/2,&\text{on }\alpha_i=a.
 \end{cases}
\]
The exact physical determinant of Lemma~\ref{lem:cs-determinant}
says $u_x\le C_{\rm side}a(v)b(v)$.

\begin{theorem}[Complete interior isoroof charge]
\label{thm:iso-interior}
For one original interior horizontal or vertical side, its
one-copy regular charge restricted to $F_v=0$, $P\ne0$ satisfies
\begin{equation}\label{eq:iso-one-side}
 \mathfrak A^{\rm side}_{i,m,R}
 \le F_{\rm cs}C_{\rm side}HE(\pi+2Q)
                         \sigma^{(L-1)/2}.
\end{equation}
The sum includes every physical word and all original exact labels
and physical roof values as positive restrictions. In particular
it bounds all regular constant-roof contributions from that side,
including all shallow indices.
\end{theorem}
\begin{proof}
Fix $\eta>0$ and split the selected source pointwise at $a\le\eta$.
On an overlap $I\cap J$ its small-$a$ charge is at most
\[
 F_{\rm cs}C_{\rm side}\min(a_I,\eta)b_J\,\pi/2
\]
by Lemma~\ref{lem:iso-balance}. The large-$a$ charge is at most
\[
 F_{\rm cs}C_{\rm side}a_Ib_J\,Q\epsilon/\eta.
\]
The cutoff $a(v)\le\eta$ need not be an interval. Both estimates
are for the complete selected subset of the overlap, because its
angular image is monotone in the one opposition quadrant.
There is no repeated payment for its connected components.

Apply the weighted interval-overlay lemma to the original two
source-disjoint partitions, first with weights
$\min(a_I,\eta)$ and $b_J$, and then with $a_I,b_J$.
The truncated family still has total weight at most $H$ and
maximum at most $\eta$. Therefore
\begin{align*}
 \sum_{I\cap J\ne\varnothing}\min(a_I,\eta)b_J
       &\le H(\epsilon+\eta),\\
 \sum_{I\cap J\ne\varnothing}a_Ib_J
       &\le H(\epsilon+A_s).
\end{align*}
These are sums of actual interval suprema; neither is inferred
from an almost-everywhere endpoint kernel. They count all actual
physical overlaps, so restricting to the stationary subsource
can only decrease the preceding positive charges.

Combining the bounds gives
\[
 \mathfrak A^{\rm side}_{i,m,R}
 \le F_{\rm cs}C_{\rm side}H
 \left\{\frac{\pi}{2}(\epsilon+\eta)
                 +\frac{Q\epsilon(\epsilon+A_s)}{\eta}\right\}.
\]
Set $\eta=\sqrt{E\epsilon}$, using $\epsilon,A_s\le E$.
Then $\epsilon\le\sqrt{E\epsilon}$ and
$\epsilon(\epsilon+A_s)/\sqrt{E\epsilon}\le2\sqrt{E\epsilon}$.
The braces are at most
$(\pi+2Q)\sqrt{E\epsilon}
=(\pi+2Q)E\sigma^{(L-1)/2}$.
This proves \eqref{eq:iso-one-side}.

The argument first applies to finite source partitions and compact
regular subsets. Exhausting positive subfamilies and increasing
compact subsets proves the full countable-source statement.
The omitted artificial partition cuts have zero $A|d\theta|$
on each compact noncritical regular chart, because $A$ and $\theta$
are smooth there; countable chart exhaustion proves the same
for their complete regular subsource. Their critical limits are
retained in Section~\ref{sec:scf-section-critical}.
All this occurs at fixed $R,m$. Its constants are uniform in the
radius, side level, exact record and roof value.
\end{proof}

\subsection{Physical singular intersections and the original side union}

\begin{lemma}[No unvisited boundary arc coincides with a canonical side]
\label{lem:iso-transverse}
At a positive-visited auxiliary finite-word germ, an unvisited-disk
tangency intersects a canonical horizontal or vertical source
transversely. It therefore cannot occupy an open canonical
section-side arc. This holds in both time directions and for every
radius in the original interval, without a genericity assumption.
\end{lemma}
\begin{proof}
Use the actual source parameter $v$ above and unfold the transverse
Jacobi field along the specified physical half-word. Up to one
fixed orientation sign, its initial displacement and angle are
\[
                 (B,D)=(Rc_i,1)\quad\hbox{or}\quad(0,1)
\]
for a horizontal or vertical source. These are the exact vectors
used in the proof of Theorem~\ref{thm:cs-half-sup}. Reversing the
physical direction changes at most that common sign.

A true free flight of length $t$ sends $(B,D)$ to $(B+tD,D)$,
and a visited positive-incidence reflection sends it to
$(B,D+2B/(Rc))$. Thus $D>0$, and the transverse displacement
is strictly positive at every positive distance after the source,
in either time direction. In particular it is positive at an
unvisited tangency foot, whose distance from the preceding actual
contact is strictly positive. Indeed the physical calculation of
Lemma~\ref{lem:xf-standoff}, now at clearance zero, gives
$t_z^2\ge(1-R_+)^2-R_+^2=3/50>0$ for an interior tangency foot.
The foot cannot be an incident contact or an exterior minimum,
since that third disk is separated from both incident disks.

For clarity, on such a free flight write its contact as $q(v)$,
its direction as $e(v)$, its signed impact parameter for the
unvisited center $z$ as
\[
 b_z(v)=(z-q(v))\cdot e(v)^\perp,\qquad
 t_z(v)=(z-q(v))\cdot e(v)>0.
\]
Differentiating gives
\[
 b_z'(v)=-q'(v)\cdot e(v)^\perp
                     -t_z(v)\,\vartheta'(v).
\]
Its absolute value is exactly the unfolded transverse displacement
$B(t_z)>0$, with the common orientation sign restored if needed.
This is the same physical reflection and impact-parameter
differential as Theorem~\ref{thm:xf-jacobian}, applied to the
two canonical initial vectors rather than a second clearance mark.
At a tangency $b_z=\pm R$, so the derivative of its actual clearance
$|b_z|-R$ is nonzero as well. This is a transverse physical
calculation, not a bulk-null argument.

At fixed word and fixed radius, these zeroes are locally isolated
on the canonical source. There are only countably many such
intersections after the full word, flight and lattice families
are included. On a positive-visited noncritical germ, $A$ and
$\theta$ extend smoothly along the auxiliary canonical source;
a singleton has zero actual $A|d\theta|$ regular trace.
Critical intersections retain their incoming fan coefficients and
are assigned to the first unvisited witness in
Corollary~\ref{cor:scf-boundary-payment}. Visited-grazing limits
use the separate trace theorem of Lemma~\ref{lem:rbv-endpoints}.
Neither type is set to zero merely because its collision-area
measure is zero.
\end{proof}

\begin{corollary}[All interior original-section regular atoms]
\label{cor:iso-all-interior}
Let $\mathfrak A^{\rm int}_{m,R}$ be the total one-copy regular
constant-roof charge on every original section side at every
interior index, with overlapping side descriptions assigned once.
For $m\ge2$ put
\begin{equation}\label{eq:iso-Bint}
 B_{\rm iso}(m)=
 10F_{\rm cs}E(\pi+2Q)
       (C_0H_{\rm hor}+\tfrac12H_{\rm ver})
       (m-1)\sigma^{(m-2)/4},
\end{equation}
and set $B_{\rm iso}(1)=0$. Then
\[
 \sup_R\mathfrak A^{\rm int}_{m,R}\le B_{\rm iso}(m),
 \qquad
 m^qB_{\rm iso}(m)\longrightarrow0
 \quad\text{for every fixed }q\ge0.
\]
For the original exact-label signed derivative and the admitted
smooth test/guard class, let $d_\ell^{\rm int}(t)$ denote these
regular atomic contributions after same-roof, same-label terms
are combined. Then
\begin{equation}\label{eq:iso-signed}
 \sum_{\ell,t}|d_\ell^{\rm int}(t)|
                   \le2a_0b_0 B_{\rm iso}(m).
\end{equation}
\end{corollary}
\begin{proof}
There are ten original horizontal and ten original vertical side
occurrences. The horizontal incidence factor is at most one.
At every interior index $L=\max(i,m-i)\ge m/2$, and hence
$\sigma^{(L-1)/2}\le\sigma^{(m-2)/4}$.
Sum Theorem~\ref{thm:iso-interior} over the actual side occurrences
and the $m-1$ indices. Lemma~\ref{lem:iso-transverse} accounts for
all possible unvisited singular intersections with these sides;
corners and artificial seams have zero regular angular-image
length, while their critical sources remain separately counted.
This gives the first inequality.

At a physical section interface the two incident exact labels
have total absolute jump at most two. The original smooth
multiplier has modulus at most $a_0b_0$ and the original guard
lies between zero and one. One must include this factor two:
the one-copy charge cannot be multiplied by two incident traces
and then counted as if their total jump were one.
Combine signed contributions at the same roof and label first,
then use the triangle inequality and the all-label bound.
This proves \eqref{eq:iso-signed}, retaining every coincidence.

Finally $0<\sigma<1$, and for each fixed $q$,
\[
 (q+1)\log m-\frac{m-2}{4}\log(59/47)\longrightarrow-\infty.
\]
This proves the claimed original-normalization payment without
changing any outer guard stage or interchanging limits.
\end{proof}

\paragraph{Scope in the complete derivative.}
Together with the separately proved endpoint horizontal
constant-roof payment and the absence of endpoint vertical
constant-roof arcs, this pays all original-section regular atoms.
Corollary~\ref{cor:scf-boundary-payment} pays their critical
boundary sources as well. No claim is made here that every
unvisited-boundary regular constant-roof arc is absent or paid.
Strictly interior critical atoms and the complete absolutely
continuous roof derivative also remain part of the original
normalized local-charge problem. The all-label arithmetic,
conditional-return and four-coordinate raw inverse targets are
unchanged.

